\section{Compatibility and uniqueness}\label{sec:compatibility}

Theorem~\ref{thm:field-isomorphism} gives an isomorphism of perfect fields
that carries the perfection of every good rational subfield to the
perfection of its match under $\Theta$. We must prove that its action on all multiplicative classes is
$\Theta$ up to a single scalar. We then determine which field
automorphisms act by scalars.

\subsection{Distinguishing divisors by rational subfields}

For a true prime divisor $v$ and a good rational subfield $E\subset K$,
consider the restriction $D_v\to A_E$. Its rank here means the dimension
of its image as a $\Lambda$-vector space.

\begin{lemma}\label{lem:restriction-rank}
Let $v$ be a true prime divisor of $K/k$ and let $E\subset K$ be a curve
subfield. The restriction $v|_E$ is trivial if and only if the image of
$D_v\to A_E$ is infinite-dimensional. If $v|_E$ is nontrivial, this image
has dimension at most one.
\end{lemma}

\begin{proof}
If $v|_E$ is trivial, reduction embeds $E$ in $Kv$. The kernel of
$V_E\to V_{Kv}$ is finite-dimensional by Lemma~\ref{lem:kummer-kernel},
whereas $V_E$ is infinite-dimensional. Dualizing and using
Lemma~\ref{lem:residue-duality} shows that the restriction of $D_v$ to
$A_E$ has infinite-dimensional image.

Otherwise $v|_E$ is a positive multiple of the order at a point $x$ of
the smooth projective curve of $E/k$. If $f$ is a unit there, its residue
is some $c\in k^\times$, and $f/c$ is a principal $v$-unit in $K$.
Every character in $D_v$ therefore kills $f$. Its restriction to $E$
factors through the value group of $v|_E$, which is cyclic. The image
has dimension at most one. This argument uses decomposition characters;
it remains valid if the ramification multiplicity is divisible by $\ell$.
\end{proof}

\begin{lemma}\label{lem:separate-divisors}
For distinct true prime divisors $v_1,v_2$ of $K/k$, there is a good
rational subfield $E\subset K$ such that $v_1|_E$ is trivial and
$v_2|_E$ is nontrivial.
\end{lemma}

\begin{proof}
Choose a normal projective model $X$ on which both valuations have
codimension-one centers $D_1,D_2$. We recall why this is possible.
For each valuation, lifts of a residue transcendence basis, together
with a uniformizer, define a rational subfield over which $K$ is finite.
Normalization of a suitable projective model of that subfield realizes
the valuation as a prime divisor. A common normal projective model
dominating the two models retains codimension-one centers: the residue
field of each new center contains that of the old center, of
transcendence degree $d-1$ over $k$.

Choose two distinct closed points $P,Q$ of $D_1\setminus D_2$, and a
smooth closed point $R$ of $X\setminus(D_1\cup D_2)$. The first choice
is possible because $d\ge2$. There is an affine open containing these
points and meeting $D_2$; write its coordinate ring as $A$, and let
$\mathfrak p_2$ be the prime ideal of $D_2$ there. Choose two independent
cotangent vectors at $R$. The Chinese remainder theorem, applied to
$\mathfrak p_2,\mathfrak m_P,\mathfrak m_Q,\mathfrak m_R^2$, produces
$a,b\in A$ satisfying
\[
\begin{array}{c|ccc}
 & D_2 & P & Q\\ \hline
 a &0&0&1\\
 b &1&1&1
\end{array}
\]
and having value zero and the prescribed independent first-order terms
at $R$. These ideals are pairwise comaximal, so the prescriptions are
independent. By Lemma~\ref{lem:good-functions}, $t=a/b$ is good.
Its residue on $D_1$ is nonconstant, since it takes different values
at $P,Q$, and its value along $D_2$ is positive. Thus $v_1$ is trivial
on $k(t)$ and $v_2$ is nontrivial there.
\end{proof}

\subsection{The scalar is global}

\begin{proposition}\label{prop:global-scalar}
Let $\Psi:V_K\to V_K$ be a compatible $\Lambda$-linear automorphism.
If $\Psi(V_E)=V_E$ for every good rational subfield $E\subset K$,
then $\Psi=a\,\mathrm{id}$ for one $a\in\Lambda^\times$.
\end{proposition}

\begin{proof}
By Proposition~\ref{prop:true-divisors}, $\Psi^*$ permutes the pairs
$I_v\subset D_v$ for true prime divisors. Since $\Psi$ fixes each
$V_E$, it preserves the rank of restriction of $D_v$ to $A_E$.
Lemmas~\ref{lem:restriction-rank} and~\ref{lem:separate-divisors}
therefore force this permutation of true divisors to be the identity.

Fix a good rational subfield $E$. Every point line of $A_E$ is the
nonzero restriction of some $I_v$. The induced dual automorphism
$\Psi_E^*=(\Psi|_{V_E})^*$ of $A_E$ therefore fixes every point line.
Write
\[
 \Psi_E^*(\ord_x)=a_x\ord_x,\qquad a_x\in\Lambda^\times,
 \quad x\in C_E(k),
\]
where $C_E$ denotes the smooth projective model of $E$. The sum of
these order characters is zero, and its coefficient families have
exactly the diagonal kernel by Lemma~\ref{lem:curve-orders}.
Continuity permits applying $\Psi_E^*$ to this sum, so all $a_x$ equal
one scalar $a_E$. Since $E$ is rational, its point orders separate
$V_E$, and hence $\Psi|_{V_E}=a_E\,\mathrm{id}$.

Let $E=k(t)$ and $F=k(s)$ be any two good rational subfields. On a
common smooth model open, $t,s$ are regular and both differentials
are nonzero. Blow up a closed point of that open. Its exceptional
divisor $v$ has
\[
 \ord_v(t-t(P))=\ord_v(s-s(P))=1.
\]
Thus the inertia line $I_v$ restricts nontrivially to a point line
of both $A_E$ and $A_F$. The scalar by which $\Psi^*$ acts on $I_v$
must be both $a_E$ and $a_F$. All these scalars agree. Good-function
classes span $V_K$ by Lemma~\ref{lem:good-functions}, proving the claim.
\end{proof}

Apply the proposition to $\Psi=\alpha_1^{-1}\Theta$, where
$\alpha:K^i\to L^i$ is supplied by Theorem~\ref{thm:field-isomorphism}.
The canonical identifications under perfection are compatible with
Milnor multiplication. Also $\alpha(E^i)$ is the perfection of the
good field matched to $E$, so $\Psi$ fixes $V_E$ for every such $E$.
Consequently $\Theta=a\alpha_1$ for a single $a\in\Lambda^\times$.

\subsection{Uniqueness up to Frobenius}

\begin{proposition}\label{prop:uniqueness}
Suppose $\beta:K^i\to K^i$ is a field automorphism preserving $k$
whose action on $V_K$ is a scalar. Then $\beta$ is the identity in
characteristic zero and an integral power of Frobenius in positive
characteristic.
\end{proposition}

\begin{proof}
Let $E=k(t)$ be good. Its perfection is relatively algebraically
closed in $K^i$: if an element is algebraic over $E^i$, clearing
Frobenius powers puts an algebraic element over $E$ in $K$, and
relative algebraic closedness then puts it in $E$. The same applies
to $\beta(E^i)$. Choose a nonconstant element of $\beta(E^i)\cap K$
by clearing powers, and let $P$ be the relative algebraic closure
of the rational field it generates in $K$. Relative algebraic closedness
of $\beta(E^i)$ gives $P^i\subset\beta(E^i)$. Conversely, every element
of $\beta(E^i)$ is algebraic over this rational field; clearing powers
puts it in $P$. Hence $\beta(E^i)=P^i$. Since $\beta_1$ is scalar, $V_P=V_E$ inside $V_K$.
Lemma~\ref{lem:curve-intersection} forces $P=E$.

It follows that $\beta(t)=h(t)$ for a perfect rational function $h$
generating $k(t)^i$ over $k$ up to perfection. This implies that $h$
belongs to the group $\mathcal H_k$ of projective transformations and
Frobenius powers from Section~\ref{sec:field}. Indeed, after clearing
powers, its separable degree must be one; a larger separable degree
could not disappear on passing to perfections. In characteristic
zero this says simply that $h$ is projective linear.

Write $\rho=\beta|_k$. For any two distinct points $a,b\in\Pj^1(k)$,
choose a rational function $f_{a,b}(t)$ with divisor $[a]-[b]$.
The point-order vector of $\beta_1([f_{a,b}])$ is supported exactly at
\[
 h^{-1}(\rho(a)),\qquad h^{-1}(\rho(b)).
\]
Here $\rho$ fixes infinity. For negative Frobenius powers, the order
vector is interpreted by clearing powers; multiplication by a power
of $p$ is invertible in $\Lambda$. Both displayed coordinates therefore
have nonzero orders. Scalar action on $V_E$ implies that the
permutation $h^{-1}\rho$ fixes every unordered pair of distinct
points. On a set with at least three points this forces every point
to be fixed: intersect the fixed pairs $\{a,b\}$ and $\{a,c\}$.
Thus $h$ and $\rho$ agree on all points.

In particular $h$ fixes $0,1,\infty$. An element of $\mathcal H_k$
with these three fixed points is a pure Frobenius power, or the
identity in characteristic zero. In positive characteristic its
exponent is determined by $\rho$ and therefore is the same for every
good $t$. Distinct Frobenius powers act differently on the infinite
algebraically closed field $k$: equality of two powers would force
every element to satisfy $X^{p^n}-X=0$ for some $n>0$.
Good functions generate $K$, so $\beta$ is this same Frobenius power
on $K$ and then on its perfection.
\end{proof}

\begin{proof}[Proof of Theorem~\ref{thm:main}]
A compatible $\Theta$ gives equal transcendence degrees by
Proposition~\ref{prop:local-dimension}. Theorem~\ref{thm:field-isomorphism}
produces a field isomorphism $\alpha:K^i\to L^i$ respecting constants;
in particular, the characteristics agree. Proposition~\ref{prop:global-scalar}
shows that $\Theta$ and $\alpha_1$ differ by a single scalar, proving
surjectivity of the stated map. If two field isomorphisms have the
same image modulo scalars, their quotient satisfies
Proposition~\ref{prop:uniqueness}, proving injectivity modulo Frobenius.
Conversely, every integral Frobenius power acts on $V_L$ by the
corresponding nonzero scalar. The map is therefore well-defined and
bijective with exactly the stated equivalence relations.
\end{proof}
